% Zdroj: import_hodiny/40_Podklad_prubezny_test.tex a jeho PDF, obě strany.
% Použity části Teorie a vysvětlení a Metodické poznámky.
% Diagnostické P1–P5 ani opravné C1–C5 nejsou převzaty jako zadání testu.
% Metodický záměr těchto úloh je zachován v přehledu dovedností a následné nápravě.
% Konkrétní numerická zadání jsou záměrně vynechána podle výslovného zadání učitele.
% Nepřebíráno bodování z importu: pro tuto přípravu nebylo požadováno.
% Rozpis: 4_rocnik/tematicke_plany/Rozpis_vyuky_MAT_4_rocnik.tex, řádek 40 (T3).
% Plán: 4_rocnik/tematicke_plany/SPS.pdf, 2.1–2.2.
% Sazba: sablony/pripravy/TEMPLATE_Priprava.tex a společný priprava.sty.
\documentclass[11pt,a4paper]{article}
% Spolecny styl pro ucitelske pripravy MAT (Work/MAT)
% Verze 1.1, 2026-10-23
\NeedsTeXFormat{LaTeX2e}
\ProvidesPackage{priprava}[2026/10/23 v1.1 Sablona ucitelskych priprav MAT]

% Zakladni typografie a cestina
\RequirePackage[T1]{fontenc}
\RequirePackage[utf8]{inputenc}
\RequirePackage[czech]{babel}
\RequirePackage{lmodern}
\RequirePackage{microtype}
\RequirePackage{amsmath,amssymb}
\RequirePackage{xcolor}
\RequirePackage{enumitem}
\RequirePackage[most]{tcolorbox}
\RequirePackage{tabularx,booktabs,array}
\RequirePackage{fancyhdr}
\RequirePackage{needspace}
\RequirePackage{geometry}
\RequirePackage[hidelinks]{hyperref}

\geometry{a4paper,top=17mm,bottom=18mm,left=23mm,right=23mm}
\definecolor{notegray}{gray}{0.39}
\setlength{\parindent}{0pt}
\setlength{\parskip}{2.5pt plus 1pt minus 0.5pt}
\setlist[itemize]{topsep=2pt,itemsep=1pt,parsep=0pt,leftmargin=1.45em}
\setlist[enumerate]{topsep=2pt,itemsep=2pt,parsep=0pt,leftmargin=1.65em}
\setcounter{secnumdepth}{0}
\renewcommand{\arraystretch}{1.12}

% Zahlavi: vlevo rocnik | globalni tema, vpravo cislo: nazev hodiny
% Zapati zustava zcela prazdne.
\pagestyle{fancy}
\fancyhf{}
\renewcommand{\headrulewidth}{0.35pt}
\renewcommand{\footrulewidth}{0pt}
\setlength{\headheight}{14pt}
\newcommand{\nastavPripravu}[4]{%
  \fancyhead[L]{\small #1\enspace\textbar\enspace #2}%
  \fancyhead[R]{\small #3: #4}%
}

% Minimalisticke cernobile boxy. Kazda definice zvlast.
\tcbset{commonlesson/.style={enhanced,breakable=false,colback=white,colframe=black,
  colbacktitle=white,coltitle=black,boxrule=0.45pt,
  arc=0.8mm,outer arc=0.8mm,left=3mm,right=3mm,top=1.6mm,bottom=1.6mm,
  toptitle=1mm,bottomtitle=0.7mm,fonttitle=\bfseries,
  attach boxed title to top left={xshift=2.5mm,yshift=-1.15mm},
  boxed title style={colback=white,colframe=white,boxrule=0pt,left=1mm,right=1mm}}}
\newtcolorbox{zadaniBox}[1][]{commonlesson,title={Zadání},#1}
\newtcolorbox{prikladBox}[1][]{commonlesson,title={Příklad},#1}
\newtcolorbox{definiceBox}[1][]{commonlesson,title={Definice},#1}

% Spolecny box pro zadani a jeho reseni, bez mezery mezi castmi.
% Pouziti:
%   \begin{prikladReseniBox}
%     Zadani ulohy.
%     \tcblower
%     \begin{align*} ... \end{align*}
%   \end{prikladReseniBox}
% Spodni cast ma jemny sedy podklad a delici linku; lze v ni sazet
% vice radku matematiky. Obe casti maji jen jeden spolecny ramecek.
\tcbset{solutionlesson/.style={
  commonlesson,
  enhanced jigsaw,
  breakable,
  colbacklower=black!3,
  segmentation style={solid,draw=black!25,line width=0.35pt},
  fontlower=\small\color{notegray},
  before lower={\noindent\textbf{Řešení:}\enspace},
  leftlower=3mm,rightlower=3mm
}}
\newtcolorbox{prikladReseniBox}[1][]{solutionlesson,title={Příklad},#1}
\newtcolorbox{zadaniReseniBox}[1][]{solutionlesson,title={Zadání},#1}

% Pomucky pro autora pripravy: sedive poznamky; \reseni pro starsi pripravy
\newcommand{\poznamka}[1]{\par\begingroup\small\color{notegray}\itshape
  \textnormal{\textbf{Poznámka pro učitele.}}\enspace #1\par\endgroup}
\newcommand{\reseni}[1]{\par\begingroup\small\color{notegray}
  \textbf{Očekávané řešení:}\enspace #1\par\endgroup}
\newcommand{\blok}[1]{\needspace{5\baselineskip}\subsection*{#1}}
\newcommand{\vyrok}[1]{\uv{#1}}

\nastavPripravu{4. ročník}{Kombinatorika}{40}{Průběžný test z kombinatoriky}
\begin{document}

Hodina je určena průběžnému testu 3 v rozsahu \textbf{variace, permutace a kombinace}. Závazný rozpis vymezuje na písemnou část \textbf{15–20 minut}; zbytek hodiny slouží krátkému rozboru a navazujícímu procvičování. Tato příprava poskytuje učiteli obsahové a metodické vedení. Konkrétní zadání ani varianty písemné práce neobsahuje.

\poznamka{Předpoklady: výuka hodin 29–39; pro písemnou část především dovednosti hodin 29–36. Cíl: zjistit, zda žák podle slovního zadání rozpozná pořadí, opakování a omezení, sestaví odpovídající početní model a vysvětlí ho. Samotná znalost názvu vzorce není dostatečným dokladem porozumění.}

\blok{Vymezení ověřovaného učiva}
Do obsahového rámce patří pravidla součtu a součinu, faktoriál jako nástroj počítání, variace bez opakování i s opakováním, permutace bez opakování i s opakováním a kombinace bez opakování. Ověřuj také rozdíl mezi číslem a kódem, nerozlišitelnost stejných kopií a jednoznačnost započtení výsledků.

Pascalův trojúhelník, samostatné důkazy identit a binomická věta nejsou podle názvu testu v rozpisu povinným zaměřením této písemné práce. Není vhodné test rozšířit na celou kombinatoriku jen proto, že se tyto části již vyučovaly. Náročné vícefázové úlohy patří především do hodiny 41.

\begin{tabularx}{\linewidth}{@{}>{\raggedright\arraybackslash}p{43mm}X@{}}
\toprule
\textbf{Ověřovaná dovednost} & \textbf{Co sledovat v žákově postupu}\\\midrule
Postupné volby a omezení & Správně přiřazené počty možností jednotlivým pozicím; respektování zákazu nuly na začátku čísla.\\
Variace bez opakování & Rozlišené role či místa; ubývající počet možností; počet vybraných osob menší nebo roven počtu dostupných.\\
Variace s opakováním & Pevná délka kódu a stálá nabídka na každém místě; uvědomění, že PIN smí začít nulou.\\
Permutace různých prvků & Použití všech prvků právě jednou; počet míst souhlasí s počtem objektů.\\
Permutace stejných kopií & Správné četnosti, jejich součet a dělení faktoriály četností.\\
Kombinace bez opakování & Výběr skupiny bez rozlišených rolí; oprava vícenásobného započtení dělením počtem pořadí.\\
\bottomrule
\end{tabularx}

\poznamka{Při samostatné přípravě budoucího testu vyber malé množství krátkých položek s různými diagnostickými cíli. Není reálné za 15–20 minut požadovat důkladnou slovní úlohu na každý uvedený typ. Konkrétní skladbu, přiřazení k variantám a čas ověř až nad vybranými úlohami; postupuje se podle skillu \texttt{vytvor-test} a výběr schvaluje učitel. Dosavadní podklad v importu není hotový schválený test.}

\newpage
\blok{Organizace 45 minut}
\poznamka{0–3 min: vysvětli organizaci a požadavek stručně zdůvodňovat postup. 3–21 min: písemná práce, tedy doporučených 18 minut v rámci rozpisu. 21–28 min: po odevzdání krátce porovnej rozhodovací strategie. 28–40 min: cílená náprava podle pozorovaných obtíží. 40–45 min: samostatná slovní rekapitulace a příprava na hodinu 41. Při jiné délce testu v rozmezí 15–20 minut uprav délku navazujícího procvičení.}

Předem připrav až následně schválené zadání a organizaci práce. Použití kalkulačky, tabulek ani dalších pomůcek tato příprava nepovoluje; podmínky stanoví učitel při zadávání konkrétního testu. Během práce upřesňuj pouze organizační pokyny, neposkytuj nápovědu k volbě vzorce jednotlivým žákům.

Po odevzdání nehodnoť před třídou individuální výkony. Krátký společný rozbor může vycházet z žákovských zdůvodnění a z pozorování učitele, ale úplné vyhodnocení proběhne až po opravě prací. Nenahrazuj ho odhadem úspěšnosti podle několika veřejných odpovědí.

\blok{Vyhodnocení a následná náprava}
Při opravování odlišuj chybný model, chybný počet možností v jednom kroku, početní chybu a nedostatečně doložený postup. Sleduj, zda žák rozumí svému správnému výsledku, i zda je schopen opravit lokální chybu při jinak správné úvaze. Tato příprava nezavádí bodovací schéma ani klasifikační hranice.

\begin{tabularx}{\linewidth}{@{}>{\raggedright\arraybackslash}p{44mm}X@{}}
\toprule
\textbf{Zjištěný problém} & \textbf{Doporučená bezprostřední náprava}\\\midrule
Komise zaměněná s funkcemi & Porovnej skupinu a obsazení rolí; nech vysvětlit, zda prohození dvou osob mění výsledek. Vrať se k hodině 35, P1.\\
Povolené opakování přehlédnuto & Porovnej stálou a ubývající nabídku znaků. Použij hodinu 34, P3, a nech popsat obě situace.\\
Počáteční nula u čísla & Nech oddělit první pozici a vyjmenovat její dovolené cifry. Navazuje hodina 30, P3.\\
Stejné kopie považovány za různé & Dočasně označ kopie a pak označení smaž; vysvětli počet shodných podob. Využij hodinu 33, P1–P2.\\
Ztráta vnitřního pořadí bloku & Ukaž obě orientace dvojice a odliš je od bloku stejných písmen. Porovnej hodinu 32, P2, a hodinu 33, P4.\\
Správný model, chyba ve výpočtu & Krátit faktoriály před vyčíslením; nechat přepočítat pouze chybný krok. Využij hodinu 31, C2 a C4.\\
\bottomrule
\end{tabularx}

Následné procvičení vybírej z již vytvořených příprav podle skutečných chyb; není nutné projít všechny řádky tabulky. Žák nejprve označí místo, kde jeho původní rozhodnutí selhalo, potom vysvětlí správný model a teprve nakonec přepočítá výsledek. Po dokončení oprav učitel sdělí jednotlivcům jejich výsledky a třídě shrne nejčastější společné problémy.

\poznamka{Závěrečné ověření může být ústní nebo krátkým zápisem: žák vyjmenuje rozhodující otázky a uvede příklad situace, v níž pořadí rozhoduje. Příští hodina 41 využije zjištění z testu při smíšených úlohách s doplňkem, blokem a rozdělením na případy. Nezačínej zde novou látku pravděpodobnosti.}
\end{document}
